\section{What do round-level detectors achieve?}
\label{app:roundmonitor}

On Qwen3-32B, whether the target leaks this round or the next is modestly predictable; the strongest text monitor scores a little above the probe on the episodes collected before the round-level plan, and about the same on the test episodes.
Whether an episode will ever leak is barely predictable by either.

\paragraph{The round-level setup.}
These analyses are separate from the ladder of \cref{fig:headline}, which uses one set of calls and one leak definition, with a label for each of its three questions.
They use the round-level label of \cref{app:labels} on the Qwen3-32B episodes scored for the round-level plan, and stop scoring before the secret is first written out, so they never include the leaking call.
Three monitor models, asked two questions at three levels of access, were scored on the same calls (\cref{app:monitors}).
The strongest is Jev asked to forecast a leak at any later point, shown everything the target had seen.

\paragraph{What a monitor should be asked and shown.}
Asked to forecast, Jev does better than asked to detect a rule break in progress.
Showing it the peer's messages, on top of the target's actions, gives most of its gain, and adding everything else the target saw adds little.
Neither the forecasting question nor more access helped the other two monitors.

\paragraph{Adding the probe to a monitor.}
A small logistic model on the round, the prompt length and a monitor's score was refit with the probe's score added.
The probe adds clearly to the weaker monitor on the test episodes, and adds nothing clear to the strongest one (\cref{tab:combination}).

\paragraph{Eventual outcome.}
Whether the episode will ever leak is near chance for the probe, and only slightly above for the strongest monitor.
Branching reruns from the same state often end differently, which suggests the outcome is often not yet decided far from the act.

\begin{table}[h]
  \centering
  \small
  \input{tables/round_level.tex}
  \caption{Round-level results on Qwen3-32B, with the round-level label (\cref{app:labels}) and its own call set, not those of \cref{fig:headline}: probe and monitors on the same calls, and the gain from adding the probe to a monitor. The development probe here is trained on the eventual label and scored on both labels without refitting.}
  \label{tab:combination}
\end{table}

\paragraph{Robustness.}
The round-level ``this round or the next'' result on Qwen3-32B holds across two label variants that use hindsight to choose negatives, a fixed middle layer, weighting scenarios equally, and a small neural network instead of logistic regression (\cref{tab:robust-round}).

\begin{table}[h]
  \centering
  \small
  \input{tables/robust_round.tex}
  \caption{Round-level robustness on Qwen3-32B, with the round-level label (\cref{app:labels}), on the same calls as \cref{tab:combination}. Each probe here is refit on the near-term label, with its layer chosen again, so it scores higher on that label than the eventual-trained probe of \cref{tab:combination}.}
  \label{tab:robust-round}
\end{table}

\subsection{Why is round-level prediction near chance on gpt-oss-20b with a channel-only label?}
\label{app:gptoss-round}

On gpt-oss-20b, with the leak defined as a channel post, the probe and the best monitor are both near chance for this round or the next.
The model usually leaks soon after the request: many of its channel leaks leave no call to score before them, and for most of the rest the earliest possible warning falls in the same round.
Under the main label, which counts the earlier file write, the ladder of \cref{fig:headline} shows the same pattern as on Qwen3-32B.
